\chapter{课程视频目录与实时追加区}

\section{视频来源}

公开课程信息：
\begin{itemize}
  \item 标题：\href{https://www.bilibili.com/video/BV1gW5u6UEML}{(已完结) 热力学与统计物理-汪志诚版}
  \item 上传账号：B站物院
  \item 当前请求对应分 P：p=2，标题为 1.1-1.4
  \item p=2 公开章节点：1.1, 1.2, 1.3, 1.4
  \item 字幕状态：公开视频 HTML 中字幕字段为 \verb|subtitle: {allow_submit:false, list:[]}|；player API 对多个分 P 返回的字幕列表也为空。
\end{itemize}

\section{公开参考资料}

本稿没有复制教材或公开视频逐字内容，主要参考公开目录和公开课件来确定覆盖范围：
\begin{itemize}
  \item 高等教育出版社《热力学·统计物理（第六版）》产品页给出的 11 章范围：热力学基本规律、均匀物质性质、相变、多元系、第三定律、不可逆过程、近独立粒子最概然分布、玻耳兹曼统计、玻色/费米统计、系综理论、涨落理论、非平衡态统计初步。
  \item 高等教育出版社《热力学·统计物理学习辅导书》产品页，用于确认该书是主教材配套习题辅导。
  \item 中国科学技术大学公开课件《热力学与统计物理》与《第九章系综理论》，用于核对本科课程常见讲授顺序和系综关键词。
\end{itemize}

\section{分 P 路线}

\begin{longtable}{c p{0.72\textwidth}}
\toprule
分 P & 标题 \\
\midrule
1 & 绪论 \\
2 & 1.1-1.4 \\
3 & 1.5-1.11 \\
4 & 1.14-1.18 \\
5 & 第一章作业题 \\
6 & 2.1 2.2 \\
7 & 2.3-2.5 \\
8 & 2.6-2.7 \\
9 & 第二章作业题 \\
10 & 3.1 3.4 \\
11 & 3.5 3.7 \\
12 & 3.8 3.9 \\
13 & 第三章作业题 \\
14 & 4.1-4.3 \\
15 & 4.4-4.8 \\
16 & 第四章作业题 \\
17 & 6.1-6.2 \\
18 & 6.3-6.5 \\
19 & 6.6-6.8 \\
20 & 第六章作业题 \\
21 & 7.1-7.2 \\
22 & 7.3-7.5 \\
23 & 7.6-7.8 \\
24 & 第七章作业题 \\
25 & 8.1-8.2 \\
26 & 8.3 \\
27 & 8.4-8.5 \\
28 & 第八章作业题 \\
29 & 9.1-9.4 \\
30 & 9.5 9.7 \\
31 & 9.9-9.11 \\
32 & 第九章作业题 \\
33 & 正则系综(李政道版) \\
\bottomrule
\end{longtable}

\section{实时追加模板}

每次看视频后，把笔记按以下格式补进对应章节：
\begin{enumerate}
  \item 视频位置：例如 p=2, 12:30；
  \item 主题：例如准静态过程与可逆过程；
  \item 我的问题：到底是符号、物理图像还是推导卡住；
  \item 现在理解：用自己的话复述；
  \item 仍不确定：后续需要查书或问老师的点。
\end{enumerate}
